% Source: https://artofproblemsolving.com/community/c3h2424379_formula_for_dividing_line_of_a_regular_hexagon
% Stable identity: aops-topic-2424379-q01
\textbf{(a)} \emph{Lines through $O$.} The half-turn about $O$ maps the hexagon onto itself and maps a line $\ell'$ through $O$ onto itself, swapping its two sides. So it maps one piece onto the other, and the two pieces have equal areas.

\emph{Lines not through $O$.} Let $\ell$ be a line not through $O$ that cuts the hexagon, and let $\ell'$ be the line through $O$ parallel to $\ell$. Let $U$ be the piece cut off by $\ell$ on the side away from $O$, and $U'$ the piece cut off by $\ell'$ on the side containing $\ell$. Then $U$ is contained in $U'$, and $U'$ also contains every point of the hexagon strictly between $\ell'$ and $\ell$, which is not in $U$. Since $O$ lies on $\ell'$ and is an interior point of the hexagon, all points strictly between the two lines and close enough to $O$ belong to the hexagon, so this part has positive area (shaded below). Hence
\[\text{area}(U)<\text{area}(U')=\dfrac12\cdot\text{area of the hexagon},\]
so $\ell$ does not bisect the area.

\begin{asy}
usepackage("amsmath");
size(4.5cm);
defaultpen(fontsize(10pt));
pair O=(0,0); pair[] V; for(int k=0;k<6;++k) V.push(dir(60k));
path H=V[0]--V[1]--V[2]--V[3]--V[4]--V[5]--cycle;
pair d=dir(20), n=dir(110); real e=0.28;
pair P1=O-1.4*d, P2=O+1.4*d, Q1=O+e*n-1.4*d, Q2=O+e*n+1.4*d;
fill(H,rgb(0.92,0.92,0.92));
pair[] a=intersectionpoints(P1--P2,H), bq=intersectionpoints(Q1--Q2,H);
fill(a[0]--a[1]--bq[1]--bq[0]--cycle,rgb(1,0.8,0.5));
draw(H,linewidth(0.8));
draw(P1--P2,blue); draw(Q1--Q2,red);
dot(O); label("$O$",O,S);
label("$\ell'$",P2,E,blue); label("$\ell$",Q2,E,red);
\end{asy}

\textbf{(b)} By (a), the bisecting lines are exactly the lines through $O$. For such a line, the half-turn about $O$ maps the segment it cuts out onto itself, so $O$ is the midpoint of the segment, and its length is $2\cdot OX$, where $X$ is the point where a ray from $O$ meets the boundary.

The boundary point $X$ lies on some side. The distance from $O$ to points of a side is smallest at the midpoint of the side, where it equals the apothem $\dfrac{\sqrt3}{2}s$ (the height of an equilateral triangle of side $s$), and largest at its endpoints, the vertices, where it equals $s$. Hence
\[\sqrt3\,s\le\text{length}\le2s.\]

\begin{asy}
usepackage("amsmath");
size(4.5cm);
defaultpen(fontsize(10pt));
pair O=(0,0); pair[] V; for(int k=0;k<6;++k) V.push(dir(60k));
draw(V[0]--V[1]--V[2]--V[3]--V[4]--V[5]--cycle,linewidth(0.8));
pair m1=(V[0]+V[1])/2, m2=(V[3]+V[4])/2;
draw(V[1]--V[4],red+linewidth(0.9)); draw(m1--m2,blue+linewidth(0.9));
pair u=dir(15); pair X1=extension(O,O+u,V[0],V[1]), X2=extension(O,O-u,V[3],V[4]);
draw(X1--X2,deepgreen+dashed);
dot(O); label("$O$",O,NW);
label("$2s$",(V[1]+O)/2,W,red); label("$s\sqrt3$",m2,W,blue);
dot(m1); dot(m2); dot(X1); dot(X2); label("$X$",X1,E,deepgreen);
label("$s$",(V[4]+V[5])/2,SE);
\end{asy}

The shortest bisecting segment, of length $\boxed{\sqrt3\,s}$, joins the midpoints of two opposite sides, and the longest, of length $\boxed{2s}$, joins two opposite vertices. Every length in between also occurs, for the lines through $O$ in between.
